\section{Explicit Evidence Boundary}
\label{sec:prof-evidence-boundary}

The manuscript does not overrun repository truth. The current repo supports a
universal runtime-safety architecture, but it does not support a flat claim
of equal empirical maturity across all six domains. The evidence boundary is
therefore part of the scientific contribution, not a defensive footnote.

\begin{table*}[t]
\centering
\caption{Allowed claims versus non-claims at the current evidence boundary.}
\label{tab:prof-evidence-boundary}
\small
\begin{tabularx}{\textwidth}{@{}p{0.24\textwidth}p{0.34\textwidth}p{0.34\textwidth}@{}}
\toprule
\textbf{Claim family} & \textbf{Supported by current repo truth} & \textbf{Not supported yet}\\
\midrule
Universal runtime grammar & One adapter contract, one benchmark schema, one CertOS lifecycle, and one degraded-observation hazard model across all rows. & A claim that all rows have identical proof depth or equal closure.\\
Battery witness depth & Deep theorem-to-code-to-artifact lineage with locked reference artifacts. & Treating battery as the conceptual center of the paper.\\
AV / industrial / healthcare & Defended bounded rows with locked replay and bounded runtime semantics. & Unqualified claims of unrestricted deployment readiness.\\
Navigation & Portability and protocol design evidence with a visible real-data gap. & Field-grade navigation validation or parity with the defended rows.\\
Aerospace & Experimental boundary evidence plus supplemental runtime diagnostics. & Proof-candidate or proof-validated aerospace closure.\\
\bottomrule
\end{tabularx}
\end{table*}

The tracked matrices make the boundary explicit. In
\texttt{reports/publication/orius\_equal\_domain\_parity\_matrix.csv}, battery is the
reference witness, autonomous vehicles, industrial process control, and healthcare
monitoring are defended bounded rows, navigation is shadow-synthetic, and aerospace is
experimental. The runtime-budget matrix and the universal validation summary say the
same thing in a different form: AV, industrial, and healthcare clear their bounded
contracts; navigation remains blocked by the missing real-data row; and aerospace
remains gated by the lack of a real multi-flight safety task with material post-repair
improvement.

The paper therefore uses the following claim discipline:
\begin{itemize}
  \item It may claim a universal safety layer for Physical AI under degraded observation.
  \item It may claim bounded validation for AV, industrial, and healthcare under their current contracts.
  \item It may claim portability and protocol design for navigation.
  \item It may claim experimental boundary evidence for aerospace.
\end{itemize}

It does not claim equal-domain universality, universal hardware closure, or field
deployment readiness for the navigation and aerospace rows.  Those would outrun the
present repository evidence and the locked validation outputs.

The practical implication is straightforward.  The manuscript can be ambitious about
architecture, but it must stay conservative about closure.  That is not a weakness in
the ORIUS program.  It is the reason the program remains defensible under serious
review.
